\section{Introduction}\label{sec-intro}

Random series priors are among the most popular choices to construct nonparametric priors in Bayesian inference. Consider the regression model given by
\begin{equation} \label{eq:model}
    y_{i} = f(x_{i}) + \varepsilon_{i}, \qquad \varepsilon_{i}\overset{\mathrm{iid}}{\sim} N(0,\sigma^{2}),
\end{equation}
for $i=1,\ldots,n$. With a given set of basis functions $\{\xi_{j}\}_{j \geq 1}$, the random series prior
induces random functions $f(x) = \sum_{j=1}^{J_n}\theta_{j} \xi_{j}(x) $ with $\theta = (\theta_1, \ldots, \theta_{J_n})$ being the regression coefficients drawn from a prior distribution $\theta \sim \Pi(\cdot)$. Here $J_n$, the number of basis functions, balances the bias--variance trade-off and is critical in driving the performance of random series, a phenomenon frequently observed in empirical studies and relatively well understood theoretically. Several strategies for selecting $J_n$ have been proposed; we describe them next.

The number of terms can be selected from the data by cross-validation \citep{Ruppert2002}, by information criteria, or by empirical Bayes. Each of these procedures commits to one truncation level, and its adaptation to unknown smoothness depends on the procedure and the setting; a choice that performs well in one application need not transfer to another.

Hierarchical Bayesian methods instead place a prior on $J_n$, which quantifies uncertainty in model complexity. Random series priors of this type attain adaptive near-minimax contraction rates \citep{Arbel2013,ShenGhosal2015}; see also \citet{deJonge2012} for tensor-product splines and \citet{GuhaHoNguyen2021} for finite mixtures with a prior on the number of components. Posterior sampling, however, must then move between spaces of different dimension.

The Bayesian shrinkage literature obtains contraction without a prior on the model dimension. For the Horseshoe in the sparse normal-means model, \citet{VanDerPas2014} established near-minimax contraction with a global scale tuned to the sparsity level, and \citet{VanDerPas2017} obtained adaptation by estimating the global scale or placing a prior on it. \citet{WeiGhosal2020} derived contraction rates for shrinkage priors, including the Horseshoe, in high-dimensional logistic regression, and \citet{Castillo2015} treat sparse linear regression; see \citet{Buhlmann2011} for a broader treatment of sparsity. Closer to the present setting, \citet{Bach2021} pair a large number of spline knots with a smoothness prior whose variance has a hyperprior, and \citet{Johnstone2005} select thresholds in a large wavelet dictionary by empirical Bayes. These results motivate a large fixed set of basis functions with shrinkage on the coefficients.

This article combines the simplicity of a fixed number of terms with the adaptivity of hierarchical approaches. We take a deterministic $J_n\asymp n^{b}$ for a constant $b>0$ and place a shrinkage prior on the coefficients; we call the result a large series prior. The contributions are twofold. First, we establish posterior contraction rates for product shrinkage priors on large dictionaries under prior-mass conditions adapted from \citet{WeiGhosal2020}. The rates are stated in the empirical prediction norm, so no compatibility condition on the design is needed and the dictionary may be redundant. Second, we show that one prior, with the number of basis functions fixed independently of the truth, attains near-minimax rates pointwise throughout its admissible smoothness range, under a sparse approximation condition that we verify for Fourier and multiresolution B-spline dictionaries and, for univariate regression functions, for the Chebyshev dictionary, the non-periodic counterpart of the Fourier dictionary, which requires no periodicity or boundary condition; for the tensor-product Chebyshev dictionary we verify the same condition, and obtain the same contraction rates, for functions on $[0,1]^d$ whose univariate sections are H\"older uniformly over the other coordinates; the Horseshoe with a deterministic global scale satisfies the prior conditions. More generally, the theory shows that the rate is governed by how sparsely the basis system represents $f_0$, and our simulations show the same dependence empirically: the observed error decay follows the sparsity of the fitted coefficients. Because the dictionary and the design are fixed, the Gram matrix of the design is computed once; for the dictionary--design pairs used in our simulations it is diagonal, diagonal plus a rank-one term, or banded, and an exact Gibbs sampler for the Horseshoe then costs $O(J_n)$ operations per sweep instead of $O(J_n^3)$.

Section~\ref{sec:method} introduces the large series prior and its assumptions, presents the posterior contraction theorem, and derives rates for H\"older-type classes; all proofs are in the Appendix. Section~\ref{sec-verify} reports simulations on coefficient recovery, on comparisons with least squares and ridge regression, and on the empirical scaling of the error, followed by three real-data applications and a description of the sampler; Section~\ref{sec-conc} concludes. The Appendix contains the proofs, the sparse approximation lemmas for the three dictionaries, and the algorithm, validation, and timings of the sampler. Theorem~\ref{thm:main} and Corollary~\ref{cor:anisotropic} are formally verified in Lean~4, with formal statements that match those in the paper. The accompanying code and the Lean~4 formalization are available in the official repository at \url{https://github.com/FbStatsQuant/Large-Prior-Official-Repo}.

\textbf{Notation.} Let $c_n$ and $d_n$ be positive sequences. We write $c_n\lesssim d_n$ if $c_n\leq C d_n$ eventually for a constant $C>0$, which coincides with the big-$O$ notation $c_n=O(d_n)$, and $c_n\asymp d_n$ if both $c_n\lesssim d_n$ and $d_n\lesssim c_n$; the little-$o$ notation $c_n=o(d_n)$ means $c_n/d_n\to0$ as $n\to\infty$, so that $c_n$ is asymptotically negligible relative to $d_n$. In particular, $c_n=o(n)$ states that $c_n$ grows strictly more slowly than the sample size. For a vector $\theta\in\mathbb{R}^{J_n}$, write $\|\theta\|_p=(\sum_j|\theta_j|^p)^{1/p}$, $\|\theta\|_\infty=\max_j|\theta_j|$, and $\operatorname{supp}\theta=\{j:\theta_j\neq0\}$; $\#S$ denotes the cardinality of a finite set $S$. For a function $g$, let
\begin{align*}
\|g\|_{p,n}=\left\{\frac1n\sum_{i=1}^n|g(x_i)|^p\right\}^{1/p}.
\end{align*}
We identify a function with its evaluation vector $(g(x_1),\ldots,g(x_n))\in\mathbb{R}^n$. All balls, projections, and coverings used below therefore take place in $\mathbb{R}^n$ equipped with the inner product $\langle u,v\rangle_n=n^{-1}\sum_i u_iv_i$, whose associated norm is $\|\cdot\|_{2,n}$.
In the theoretical results, the design points $x_1,\ldots,x_n$ are deterministic and the noise variance $\sigma^2$ is known; Section~\ref{sec-verify} states how each computation treats $\sigma$. Write $Y^{(n)}=(y_1,\ldots,y_n)$, and let $P_0$ and $E_0$ denote probability and expectation under model~\eqref{eq:model} with regression function $f_0$. For an evaluation vector $v\in\mathbb{R}^n$, let $P_v$ and $E_v$ denote probability and expectation when the observations are independent $N(v_i,\sigma^2)$. If $v$ is the evaluation vector of a function $f$, we also write $P_f$ and $E_f$. The symbols $p_\theta$ and $p_0$ denote the corresponding joint densities under $f_\theta$ and $f_0$, respectively.
